%% biblatex-names-numeric-test - biber's name filters under acmnumeric.
\documentclass[acmsmall,nonacm,natbib=false]{acmart}
\pagestyle{empty}
\settopmatter{printacmref=false}
\RequirePackage[
  datamodel=acmdatamodel,
  style=acmnumeric,
  ]{biblatex}
\addbibresource{biblatex-names-test.bib}
\begin{document}
\thispagestyle{empty}

\section*{Sorting}
\cite{NsAbe,NsAeZed,NsDeZed,NsAlHakim,NsFox}: a short dash-joined prefix leaves a name part before biber compares it, so the part files under the stem behind it whatever the case of the prefix, and a character command is one letter short of that pattern, so it files where its character does.

\cite{NsIbnSina}: three letters are too many for the pattern.

\cite{NsAbe,NsProt,NsFox}: a dash the braces protect is not the dash the pattern looks for, so the prefix stays and files with the rest of the part.

\cite{NsTitle}: a title is not a name at all.

\section*{Initialling}
Initials are read off the name as written, after the same prefix goes -- but only a lowercase one \citeauthor{NiNoinit} \citeauthor{NiQuirk} \citeauthor{NiUpper} \citeauthor{NiRho}.
A braced given name is one word \citeauthor{NiBraced} \citeauthor{NiSage}, a hyphen protected by braces does not split it \citeauthor{NiProtect} \citeauthor{NiTell}, and a plain hyphen does \citeauthor{NiHyphen} \citeauthor{NiVane}.
A character command is one letter here too \citeauthor{NiChar} \citeauthor{NiZeta}, and so is an accent, whatever braces protect it \citeauthor{NiAccent} \citeauthor{NiYew}.
An initial that opens with a diacritic takes the letter behind it as well \citeauthor{NiAli} \citeauthor{NiWard}.
Two such marks are two characters to biber and one quote once typeset, and a third is left out of the count \citeauthor{NiAliTwo} \citeauthor{NiXuThree} \citeauthor{NiXu}.
An initial of nothing but that quote keeps its period in an ordinary citation too \cite{NiAliTwo,NiAli,NiWard}.
It keeps that period only where nothing precedes it, and loses it behind an earlier entry \cite{NsAbe,NiAliTwo}.
A label that merely opens with a period of its own keeps it wherever it sits \cite{NzBang,NzDotNet}.

\section*{Spelling}
\cite{NzEsc,NzLit}: an accent spelled as a command and the same accent typed whole are one author, not two.

\cite{NzDotless,NzDotted}: an accent over a dotless \textbackslash{}i is that same author too.

\printbibliography

\end{document}
